
\subsubsection{Floyd-Warshall (Todos os Pares)}
O algoritmo de Floyd-Warshall resolve o problema do caminho mínimo entre
\textbf{todos os pares de vértices} do grafo simultaneamente. Ele utiliza
Programação Dinâmica, avaliando se o caminho de $i$ para $j$ passando por um
vértice intermediário $k$ é menor que o caminho direto conhecido. A
complexidade de tempo é $\mathcal{O}(V^3)$.

\paragraph{Pseudo-Código: }\mbox{}\\
\begin{algorithm}[H]
    \SetAlgoLined
    \SetKwFunction{FloydWarshall}{FloydWarshall}
    \SetKwProg{Fn}{Função}{:}{}

    \Fn{\FloydWarshall{$matrizAdj, V$}}{
        $dist \gets matrizAdj$\;

        \Para{$k \gets 0$ \Ate $V - 1$}{
            \Para{$i \gets 0$ \Ate $V - 1$}{
                \Para{$j \gets 0$ \Ate $V - 1$}{
                    \Se{$dist[i][k] + dist[k][j] < dist[i][j]$}{
                        $dist[i][j] \gets dist[i][k] + dist[k][j]$\;
                    }
                }
            }
        }
    }
    \caption{Algoritmo de Floyd-Warshall}
\end{algorithm}

\paragraph{Implementação: }\mbox{}\\
\textbf{Python: }
\begin{lstlisting}[language=Python]
def floyd_warshall(n, adj_matrix):
    # Copia a matriz de adjacencia
    dist = [row[:] for row in adj_matrix]
    
    for k in range(n):
        for i in range(n):
            for j in range(n):
                if dist[i][k] + dist[k][j] < dist[i][j]:
                    dist[i][j] = dist[i][k] + dist[k][j]
                    
    return dist
\end{lstlisting}

\vspace{0.3cm}
\textbf{C++:}
\begin{lstlisting}[language=C++]
void floyd_warshall(int n, vector<vector<int>>& dist) {
    // dist ja eh inicializada com a matriz de adjacencia
    for (int k = 0; k < n; ++k) {
        for (int i = 0; i < n; ++i) {
            for (int j = 0; j < n; ++j) {
                if (dist[i][k] < 1e9 && dist[k][j] < 1e9) {
                    dist[i][j] = min(dist[i][j], dist[i][k] + dist[k][j]);
                }
            }
        }
    }
}
\end{lstlisting}